\section{funzioni di più variabili}

Uno degli argomenti principali del secondo corso di Analisi Matematica è lo studio 
delle funzioni di più variabili. 
Il nostro ambiente di studio privilegiato sarà 
lo spazio vettoriale reale di dimensione $n$. 
Si considera assimilata la conoscenza degli spazi vettoriali di dimensione 
finita. 
In particolare si considera come fatto noto che ogni spazio 
vettoriale reale di dimensione $n$ è isomorfo a $\RR^n$. 
Spesso lo spazio vettoriale sarà naturalmente dotato di prodotto scalare (spazio euclideo) e,
anche in questo caso, sappiamo che è possibile identificare il prodotto scalare 
dello spazio euclideo con il prodotto scalare canonico di $\RR^n$.
Dal prodotto scalare si ottiene la norma euclidea, che induce una metrica e quindi una topologia.

Riprendiamo brevemente questi concetti.
\begin{definition}
L'insieme $\RR^n$ è l'insieme delle $n$-uple (ordinate) di numeri reali 
$ x = (x_1,\dots,x_n)$.
Le operazioni di addizione e di moltiplicazione per uno scalare $t\in \RR$
sono definite come segue:
\begin{align*}
 x +  y &= (x_1+y_1,\dots,x_n+y_n)\\
t  x &= (t x_1,\dots,t x_n).
\end{align*}
Con queste operazioni $\RR^n$ risulta essere uno spazio vettoriale reale di dimensione $n$.

Possiamo definire il prodotto scalare canonico di $\RR^n$ come segue:
\[
 x \cdot  y = x_1 y_1 + \dots + x_n y_n.
\]
La norma euclidea di $ x$ è definita come
\[
\abs{ x} = \sqrt{ x\cdot  x} = \sqrt{x_1^2 + \dots + x_n^2}.
\]
\end{definition}

Quando $n=2$ o $n=3$ può essere comodo (e certamente è usuale) nominare 
le coordinate di $\RR^2$ come $x$ e $y$ e le coordinate di $\RR^3$ come $x$, $y$ e $z$
(invece che $(x_1,x_2)$ e $(x_1,x_2,x_3)$).

Considereremo funzioni con dominio in $A\subset \RR^n$ e codominio $\RR^m$. 
Scriveremo $ f\colon A \subset \RR^n \to \RR^m$.
Avremo quindi $ y =  f( x)$ con $ x =(x_1,\dots,x_n)\in A$ e $ y = (y_1,\dots,y_m)\in \RR^m$.
Nel caso $m=1$ la funzione si dice scalare, 
mentre nel caso $m>1$ la funzione si dice vettoriale.
Nel caso $n=1$ la funzione si dice di una variabile, 
mentre nel caso $n>1$ la funzione si dice di 
più variabili ($n$ variabili per la precisione).
Una funzione vettoriale $ f$ può essere identificata con una 
$m$-upla di funzioni scalari $f_1,\dots,f_m$ definite su $A$:
\[
   f(x_1, \dots, x_n) = \begin{pmatrix}
    f_1(x_1,\dots, x_n) \\
    \vdots \\
    f_m(x_1,\dots, x_n)
  \end{pmatrix}.
\]

Se $I\subset \RR$ è un intervallo, 
una funzione $ u\colon I\subset \RR \to \RR^n$ si chiama 
\emph{curva}. 
Se pensiamo a $t\in I$ come il \emph{tempo}, la funzione 
$t\mapsto  u(t)$ rappresenta la dinamica di un punto 
che si muove nello spazio $\RR^n$.
L'immagine della funzione $ u(I) = \ENCLOSE{ u(t)\colon t\in I}$
si chiama \emph{traiettoria} del punto $ u(t)$ o \emph{supporto}
della curva $ u$.
A volte si denota con $[ u]$.

Se ad esempio fissiamo un punto $ x_0\in \RR^n$ 
e un vettore $ v\in \RR^n$, la curva $ u(t) =  x_0 + t  v$
rappresenta il moto rettilineo uniforme di un punto che al tempo $t=0$ 
si trova in posizione $x_0$ e che si muove con velocità $ v$.
Il supporto di $ u$ è in questo caso una retta passante per $x_0$ 
con direttrice $ v$.

Per ricondursi allo studio di funzioni di una variabile 
è spesso utile studiare la composizione della funzione di più variabili 
con opportune curve. 
Se $ f\colon A \subset \RR^n\to \RR^m$ e $ u\colon I\subset \RR\to A\subset \RR^n$ 
possiamo considerare la funzione $ g\colon I \to \RR^m$ 
data da $ g(t) =  f( u(t))$. 
Si avrà $ g = (g_1,\dots, g_m)$ con $g_k(t) = f_k( u(t))$
dove $ f = (f_1,\dots, f_m)$. 
In definitiva le funzioni $g_k\colon I\subset \RR \to \RR$ 
sono funzioni scalari di una variabile, e le possiamo studiare 
come si è fatto nel corso di Analisi Matematica Uno.

\begin{example}[funzioni lineari]
Fissato $ w \in \RR^n$ consideriamo la funzione 
lineare $f \colon \RR^n \to \RR$ definita da 
\[
   f( x) = ( x ,  w).
\]  
Fissato $ q\in \RR^n$ e $ v \in \RR^n$ consideriamo la retta (parametrizzata da)
$ u(t) = t  v +  q$.
Si avrà allora 
\[
  f( u(t)) = (t  v +  q,  w) = t ( v, w) + ( q,  w).
\]
Si osserva ad esempio che se $ v$ è ortogonale a $ w$,
cioè $( v,  w)=0$, si ha $f( u(t))= ( q,  v)$ è costante.
Abbiamo trovato una \emph{curva di livello}, cioè una curva su 
cui $f$ risulta costante.
\end{example}


\section{continuità e limiti}
La norma euclidea definisce una metrica e quindi una topologia 
su $\RR^n$. 

\begin{definition}[topologia]
Dato $ x \in \RR^n$ e $\rho>0$, la palla aperta di raggio $\rho$ centrata in $ x$ è
\[
B_\rho( x) = \{ y \in \RR^n : \abs{ y -  x} < \rho\}.
\]

Se esiste $\rho>0$ per cui $B_\rho( x) \subset A$, diremo che $ x$ è un punto interno di $A$
e diremo che $A$ è un intorno di $ x$.
Possiamo dunque definire la famiglia di intorni di $ x$ come segue:
\[
  \mathcal U_x = \ENCLOSE{U\subset \RR^n\colon \exists \rho>0\colon B_\rho( x) \subset U}.
\]

La \emph{parte interna} di $A$ è l'insieme dei punti interni di $A$.
Un sottoinsieme $A\subset \RR^n$ si dice aperto se 
è intorno di ogni suo punto. 

Un sottoinsieme $A\subset \RR^n$ si dice chiuso se il suo complementare $\RR^n\setminus A$ è aperto.

Un punto $ x\in \RR^n$ si dice essere un punto di accumulazione per $A$ 
se ogni intorno di $ x$ contiene almeno un punto di $A$ diverso da $ x$.
\end{definition}

Una successione $ a_k\in \RR^n$ è una funzione $ a \colon \NN\to \RR^n$.
Ogni punto $ a_k$ è una $n$-upla: $ a_k = ((a_k)_1,\dots, (a_k)_n) \in \RR^n$.

Una successione $ a_k$ converge ad un punto $ x\in \RR^n$ 
se $\abs{ a_k -  x}\to 0$.
Scriveremo in tal caso $ a_k \to  x$ (per $k\to +\infty$).
Notiamo che $\abs{ a_k- x}$ è una successione di numeri reali 
non negativi, e dunque l'abbiamo già studiata nel corso di \emph{Analisi Matematica Uno}.
In particolare, se espandiamo la definizione di limite, si vede 
che $ a_k \to  x$ significa:
\[
 \forall \eps>0\, \exists N\in \NN\, \forall k\ge N\colon \abs{ a_k-  x} < \eps.
\]

\begin{definition}[continuità]
Sia $f\colon A \subset \RR^n \to \RR^m$.
Fissato $x_0 \in A$ diremo che $f$ è \emph{continua} nel punto 
$x_0$
se 
\[
\forall \eps>0\, \exists \delta>0 \, \forall x\in A\colon
  \abs{x - x_0} < \delta \implies \abs{f(x) - f(x_0)} < \eps.
\] 

Una funzione si dice essere \emph{sequenzialmente continua} nel punto 
$x_0$ se per ogni successione $a_k\in A$ con $a_k\to x_0$
si ha $f(a_k)\to f(x_0)$.

Una funzione $f\colon A \subset \RR^n \to \RR^m$ si dice essere 
continua (o sequenzialmente continua) se è continua (o sequenzialmente 
continua) in ogni punto $x\in A$.
\end{definition}

\begin{theorem}[equivalenza tra continuità e continuità sequenziale]
Sia $f\colon A\subset \RR^n\to \RR^m$ e sia $x_0\in A$.
Allora $f$ è continua in $x_0$ se e solo se $f$ è sequenzialmente 
continua in $x_0$.
\end{theorem}
\begin{proof}
La dimostrazione è la stessa già vista nel corso di Analisi Matematica Uno.
\end{proof}

\begin{theorem}[continuità della funzione composta]
Sia $f\colon A\subset \RR^n \to B \subset \RR^k$ una funzione 
continua (nel punto $x_0\in A$) e sia $g\colon B\subset \RR^k \to \RR^m$ 
una funzione continua (nel punto $y_0=f(x_0)\in B$).
Allora la composizione $g\circ f \colon A \subset \RR^n \to \RR^m$ 
è continua (nel punto $x_0$).
\end{theorem}
%
\begin{proof}
La dimostrazione è la stessa già vista nel corso di Analisi Matematica Uno.
\end{proof}



\section{differenziale}

Se $V$ e $W$ sono spazi vettoriali
denotiamo con $\L(V,W)$ lo spazio vettoriale delle funzioni lineari 
e continue da $V$ in $W$.
Se $V$ ha dimensione finita non serve richiedere la continuità, 
che è sempre garantita dal seguente.

\begin{theorem}[continuità delle applicazioni lineari]
  \label{th:lineare_continua}
Ogni applicazione lineare $L\colon \RR^n\to \RR^m$ 
è continua.
\end{theorem}
%
\begin{proof}
Sia $ e_1,\dots,  e_n$ la base canonica di $\RR^n$.
Consideriamo
\[
  M = \max\ENCLOSE{\abs{L[ e_1]}, \dots, \abs{[L e_n]}}.
\]
Si ha allora, per ogni $ x\in \RR^n$,
\[
  \abs{L[ x]}
  =\abs{L[x_1  e_1 + \dots + x_n  e_n]}
  \le \abs{x_1 L[ e_1]} + \dots + \abs{x_n L[ e_n]}
  \le nM\abs{ x}.
\]
Dunque per $ x \to 0$ si ha $\abs{L[ x]}\to 0$.
Significa che $L$ è continua in $0$.

Ma in ogni altro punto si ha 
\[
  L[ x] - L[ x_0] = L[ x -  x_0] 
\]
e dunque se $ x\to  x_0$ si ottiene che $L[ x]\to L[ x_0]$.
\end{proof}

\begin{definition}[norma operatoriale]
Sia $L\colon \RR^n \to \RR^m$ una applicazione lineare.
Definiamo la \emph{norma} di $L$ come:
  \[
  \Abs{L} = \sup \ENCLOSE{\abs{L[ x]}\colon \abs{ x}\le 1}.
  \]

Si noti che essendo $L[v]-L[w] = L[v-w]$ per linearità di $L$, 
la costante $\Abs L$ non è altro che la costante di Lipschitz di $L$.
\end{definition}

\begin{theorem}[proprietà della norma operatoriale]
  Sia $L\colon V \to W$ una applicazione lineare. 
  Allora 
  \begin{enumerate}
    \item $\Abs L < +\infty$;
    \item $\Abs {tL} = \abs{t}\cdot \Abs{L}$ per ogni $t\in \RR$;
    \item $\Abs{L+M} \le \Abs L + \Abs M$ per ogni $M\in \L(\RR^n,\RR^m)$;
    \item se $\Abs L = 0$ allora $L=0$; 
    \item $\Abs L$ è la costante ottimale per cui vale 
    la stima:
    \begin{equation}
      \label{eq:stima_norma_operatoriale}  
      \abs{L[ v]} \le \Abs L \cdot \abs { v},
      \qquad \forall  v \in \RR^n,
    \end{equation}

    inoltre esiste $ v\in\RR^n$ per cui vale 
    l'uguaglianza;
    \item se $M\in \L(W,U)$ è un'altra applicazinoe 
    lineare si ha 
    \[
      \Abs{M\circ L} \le \Abs M \cdot \Abs L.
    \]
  \end{enumerate}
\end{theorem}
%
\begin{proof}
Per il teorema~\ref{th:lineare_continua} la funzione $\abs{L}$ è continua,
dunque per il teorema di Weiestrass ammette massimo sul disco
unitario $\ENCLOSE{ x \in \RR^n\colon \abs{x}\le 1}$.
Dunque $\Abs{L}\in \RR$ è finito ed è realizzato 
su un vettore di norma unitaria.
Se $t\in \RR$ e $ v$ è unitario si ha
\[
  \abs{L[t v]} = \abs{t}\cdot \abs{L[ v]}
  \le \abs{t}\cdot \Abs L \cdot \abs{ v} 
  = \Abs L \cdot \abs{t v}
\]
da cui si ottiene anche~\ref{eq:stima_norma_operatoriale}.
Ovviamente se $\Abs L=0$ si ottiene $L[v]=0$ per ogni $v\in V$ 
e dunque $L=0$.
La condizione $\Abs{L+M}\le \Abs{L} + \Abs{M}$ 
discende direttamente dalle proprietà del $\sup$, ricordando 
che per definizione $(L+M)[v] = L[v] + M[v]$.
Per l'ultimo punto osserviamo che per ogni $ v\in \RR^n$ si ha 
\[
  \Abs{ML[ v]} = \Abs{M[L[ v]]}
  \le \Abs{M}\cdot \abs{L[ v]}
  \le \Abs{M}\cdot \Abs{L}\abs{ v}.
\]
Passando al $\sup$ sui $ v$ unitari si ottiene il risultato.
\end{proof}

\begin{definition}[differenziale di Fréchet]
  Siano $V$ e $W$ spazi vettoriali normati.
  Sia $f\colon A \subset V \to W$ e sia $x_0\in A$.
  Diremo che $f$ è differenziabile in $x_0$ se esiste 
  $L\in \L(V,W)$ tale che 
  \begin{equation}\label{eq:def_differenziale}
    f(x) = f(x_0) + L[x- x_0] + o(\abs{x - x_0}), 
    \qquad \text{per } x\to x_0.
  \end{equation}
  L'applicazione lineare $L$
  si chiama \emph{differenziale} di $f$ nel punto $x_0$
  e si denota con $Df(x_0)$.
\end{definition}

\begin{remark}
Le notazioni per le derivate delle funzioni di più variabili cambiano 
a seconda degli ambiti di applicazione in maniera spesso incoerente.
Invece che $Df(x_0)$ si può usare il simbolo $\nabla f(x_0)$ oppure,
nell'ambito della geometria differenziale, il simbolo $df_{x_0}$. 
Cercheremo in questo testo di mantenere una notazione coerente, ma
invitiamo il lettore a comprendere che testi diversi possono avere 
notazioni diverse. 
\end{remark}

\begin{theorem}[proprietà delle funzioni differenziabili]
Sia $f\colon A\subset V \to W$ una funzione 
differenziabile nel punto $x_0\in A$. Allora
\begin{enumerate}
  \item $f$ è continua in $x_0$;
  \item se $x_0$ è un punto interno di $A$
  il differenziale è unico
\end{enumerate}
\end{theorem}
\begin{proof}
Visto che le applicazioni lineari sono continue
si ha $Df(x_0)[x -  x_0]\to 0$ per $x \to x_0$
e dunque dalla definizione~\eqref{eq:def_differenziale} 
si ottiene che $f(x) - f(x_0)\to 0$
per $x \to x_0$. 

Siano ora $L,M\in \L(V,W)$ che soddisfano entrambi la condizione~\eqref{eq:def_differenziale}
e sia $v\in V$ un qualunque vettore unitario.
Facendo la differenza si ottiene
\[
  (L-M)[x - x_0] = o(\abs{x-x_0}), 
  \qquad \text{per } x \to x_0.
\]
Preso $x =  x_0 + t v$ osserviamo che essendo $ x_0$ interno 
ad $A$ avremo che per $t$ sufficientemente piccolo si ha effettivamente 
$ x \in A$. 
Dunque, per $t\to 0$ si ha $ x \to  x_0$ e quindi
\[
  (L-M)[t  v] = o(\abs{t}\cdot \abs{ v}) = o(t) = t o(1)
\] 
da cui 
\[
  (L-M)[ v] = o(1).
\]
Visto che il lato sinistro non dipende da $t$ 
significa che $(L-M)[ v]=0$ per ogni $ v$ unitario.
Questo significa che $L-M = 0$ ovvero $L=M$.
\end{proof}

Se $V$ ha dimensione finita, scelta una base $e_1,\dots, e_n$ di $V$ 
è possibile mettere in corrispondenza $V$ con $\RR^n$ tramite 
un isomorfismo che associa ad ogni $v\in V$ la $n$-upla delle sue coordinate $x=(x_1,\dots,x_n)$ 
rispetto alla base scelta.
Per questo non è riduttivo supporre che dominio e codominio siano
$V=\RR^n$ e $W=\RR^m$, nel qual caso le coordinate coincidono con i vettori.
Se $f\colon \RR^n \to \RR^m$ per ogni $x\in \RR^n$ l'applicazione lineare 
$Df(x)\colon \RR^n \to \RR^m$ 
si identifica con una matrice $m\times n$
che con lieve abuso di notazione denoteremo sempre con $Df(x_0)$
(potremmo usare la notazione $df_{x_0}$ per denotare specificatamente 
l'applicazione lineare).
Dunque $Df(x_0)[v] = Df(x_0)\, v$ si può intendere come prodotto matrice/vettore.
La matrice $Df$ si chiama \emph{matrice Jacobiana} 
e a volte si denota con $J_f( x_0)$.

Nel caso $m=1$, cioè $f\colon \RR^n \to \RR$, la matrice Jacobiana
è formata da una singola riga e si identifica quindi con un 
vettore di $\RR^n$. 
Tale vettore si chiama \emph{gradiente} di $f$ e si indica con 
$\nabla f(x)$. 
Si ha dunque
\[
  Df(x)\, v = \langle \nabla f(x), v\rangle.
\]

In generale se $f\colon \RR^n \to \RR^m$ 
si può scrivere $f = (f_1,\dots,f_m)$ con $f_k\colon \RR^n \to \RR$
e si troverà quindi che le righe della matrice $Df$ jacobiana di $f$
sono i vettori gradiente $\nabla f_1, \dots, \nabla f_m$.

Nel caso $n=1$, cioè $ f\colon A\subset \RR\to\RR^m$, $x\in A$,
la matrice Jacobiana è formata da una singola colonna che si denota con
$f'(x) \in \RR^n$.
Il differenziale si scrive nella forma $Df(x)[h] = h f'(x_0)$.
La definizione di differenziale diventa
\[
  f(x) =  f(x_0) +  f'(x_0)\cdot (x-x_0) + o(x-x_0)
\] 
che è equivalente a 
\[
  \frac{ f(x) -  f(x_0)}{x-x_0} =  f'(x_0) + o(1)
\]
cioè la definizione di derivata vista nel corso di Analisi Matematica Uno.
Per le funzioni di una variabile \emph{differenziabile} è quindi equivalente 
a \emph{derivabile}.

\begin{example}\label{ex:48993}
  Si consideri la funzione $f\colon \RR^2\to \RR$, $f(x,y)=(x-1)(y+2)$.
  Si ha 
  \[
  f(x,y) = -2 + 2x - y + xy
   = f(0,0) + \langle(2,-1),(x,y)\rangle + xy.
  \]
  Visto che $\abs{xy} \le x^2+y^2 = o(\sqrt{x^2+y^2})$
  abbiamo verificato che $f$ è differenziabile nel punto $(0,0)$ 
  e che $\nabla f(0,0) = (2,-1)$.
\end{example}

\begin{theorem}[differenziale della composizione]
  Sia $ f\colon A \subset \RR^n \to B\subset \RR^k$,
  $ g\colon B\subset \RR^k \to \RR^m$, $x_0\in A$.
  Se $ f$ è differenziabile in $x_0$ 
  e $ g$ è differenziabile in $y_0 =  f(x_0)$ 
  allora la funzione composta $g\circ  f\colon A \subset \RR^n \to \RR^m$
  è differenziabile in $x_0$ e si ha 
  \[
    D (g \circ  f)(x_0) 
    = D g(y_0) D f(x_0).
  \]
  Dove il prodotto a secondo membro indica la composizione 
  delle applicazioni lineari ovvero il prodotto delle matrici corrispondenti.
\end{theorem}
%
\begin{proof}
  Facendo il cambio di variabile $y = f(x)$ nella formula~\eqref{eq:def_differenziale}
  applicata a $g$ nel punto $y_0$
  e inserendo la formula corrispondente per $f$ nel punto $x_0$ si ottiene
  \begin{align*}
     g(f(x)) 
    & =  g( y_0) 
        + D g(y_0)[ f(x)- f(x_0)]
        + o(\abs{f(x)-  f(x_0)})\\
    & =  g(y_0) 
        + D g(x_0)[D f(x_0)[x - x_0] 
          + o(\abs{ x - x_0})]
        + o(\abs{D f(x_0)[x - x_0] 
          + o(\abs{x - x_0})})\\
    & =  g(f(x_0))
       + Dg(x_0)[Df(x_0)[x - x_0]]
       + (\Abs{Dg(x_0} + \Abs{Df(x_0)} + 1)\cdot o(\abs{ x -  x_0})
  \end{align*}
  che è la condizione di differenziabilità per $ g \circ  f$ 
  nel punto $ x_0$ con differenziale $Dg(y_0) Df(x_0)$.
\end{proof}

\subsection{derivata direzionale}

Nel caso particolare in cui 
$f\colon A\subset \RR^n \to \RR^m$ e $u\colon I\subset \RR\to A\subset \RR^n$,
la formula per la derivata della funzione composta diventa
\[
  \frac{d}{dt}  f(u(t))
  = Df(u(t))[u'(t)]
\]
che nel caso $m=1$ si può anche scrivere usando il gradiente
\[
  \frac{d}{dt} f(u(t))
  = \langle \nabla f( u(t)), u'(t)\rangle.
\]

Dato $v\in\RR^n$ possiamo considerare la retta con direzione $v$
passante per un punto $x\in \RR^n$: $u(t) =  x + t v$.
Definiamo allora la \emph{derivata direzionale} di $f$ nel punto $x$
in direzione $v$ come la derivata della funzione composta:
\[
   \frac{\partial  f}{\partial  v}(x)
   = \Enclose{\frac {d}{dt}  f(u(t))}_{t=0}.
\]
E' possibile che la derivata direzionale esista anche se la funzione $f$ 
non è differenziabile, perché questa definizione coinvolge i valori della funzione 
solamente sulla retta scelta.
Ma se $f$ è differenziabile in $x$, visto che $u(t)= x + t v$ 
è derivabile con derivata $u'(t) =  v$ si ha 
\[
  \frac{\partial f}{\partial v} (x)
  = Df(x) [v].
\]
E nel caso scalare $m=1$:
\[
  \frac{\partial f}{\partial  v} ( x)
  = \langle \nabla f(x),  v\rangle.
\]

\begin{example}
Si consideri la funzione $f(x,y) = (x-1)(y+2)$.
Scelto $v=(\alpha,\beta)\in \RR^2$ possiamo calcolare
\[
\frac{\partial f}{\partial v}(0,0)
= \frac{d}{dt}\Enclose{(\alpha t-1)(\beta t+2)}_{t=0}
= \frac{d}{dt}\Enclose{-2 + (2\alpha-\beta)t + \alpha\beta t^2}_{t=0}
= \Enclose{2\alpha-\beta + 2\alpha\beta t}_{t=0}
= 2\alpha-\beta.
\]
Effettivamente sapendo dall'esempio~\ref{ex:48993}
che $\nabla f(0,0) = (2,-1)$ si verifica che 
\[
  2\alpha - \beta = \langle \nabla f(0,0),(\alpha,\beta)\rangle. 
\]
\end{example}

\begin{example}
Si consideri la funzione $f\colon \RR^2\to\RR$ definita da 
\[
  f(x,y) = \begin{cases}
    \frac{x^2 y}{x^2+y^2} &\text{se } (x,y)\neq (0,0)\\
    0 &\text{se } (x,y)=(0,0).
  \end{cases}
\]
Se fissiamo un qualunque un vettore $v=(\alpha,\beta)$ possiamo facilmente 
calcolare la derivata direzionale nel punto $(0,0)$:
\[
  \frac{\partial f}{\partial v}(0,0) 
  = \frac{d}{dt}\Enclose{f(\alpha t,\beta t)}_{t=0}
  = \frac{d}{dt}\Enclose{\frac{\alpha^2\beta}{\alpha^2+\beta^2} t}_{t=0}
  = \frac{\alpha^2 \beta}{\alpha^2+\beta^2}.
\]
Osserviamo che in questo caso $\frac{\partial f}{\partial v}$ non 
è lineare in $v$ e quindi $f$ non può essere differenziabile in $(0,0)$.
\end{example}


\subsection{derivate parziali}

Si capisce che per determinare il differenziale di una funzione 
differenziabile è sufficiente 
calcolare le derivate direzionali lungo gli $n$ vettori di una base.
In particolare se si scelgono i vettori della base canonica si ottengono
le cosiddette \emph{derivate parziali}:
\[
   \frac{\partial  f}{\partial x_k} 
   = \frac{\partial  f}{\partial  e_k}
   = \Enclose{\frac{d}{dt}  f( x + t e_k)}_{t=0}
   = \frac{d}{dx_k} f(x_1, \dots, x_k, \dots, x_n)
\]
Calcolare le derivate parziali è molto semplice perché 
si tratta di derivare la funzione $x_k \mapsto f(x_1,\dots,x_n)$
come se fosse una funzione di una sola variabile, $x_k$, mentre 
le altre variabili $x_j$ con $j\neq k$ vanno considerate costanti.

Se poi scriviamo $f = (f_1,\dots,f_m)$ possiamo quindi 
determinare tutte le componenti della matrice jacobiana:
\[
  \enclose{D f( x)}_{jk} = D f_j( x)\,  e_k
  = \langle \nabla f_j( x), e_k\rangle
  = (d f_j)_{ x}[ e_k]
  = \frac{\partial f_j}{\partial x_k}( x).
\]
Dunque 
\[
  D f( x) = \begin{pmatrix}
    \frac{\partial f_1}{\partial x_1} & \dots & \frac{\partial f_1}{\partial x_n}\\
    \vdots & \ddots & \vdots \\
    \frac{\partial f_m}{\partial x_1} & \dots & \frac{\partial f_m}{\partial x_n}\\
  \end{pmatrix}
\]
e nel caso $m=1$ 
\[
  \nabla f ( x) = \begin{pmatrix}
    \frac{\partial f}{\partial x_1}\\
    \vdots\\
    \frac{\partial f}{\partial x_n}
  \end{pmatrix}.
\]

Useremo anche una notazione più compatta che utilizza 
gli indici per le derivate parziali:
\[
   f_{x_k} = \frac{\partial f}{\partial x_k}.
\]

Le derivate parziali sono uno strumento molto 
efficace per calcolare le componenti del differenziale 
di una funzione di più variabili.
Il problema è che l'esistenza delle derivate parziali 
non garantisce la differenziabilità.
Risulta quindi molto utile il seguente.

\begin{theorem}[del differenziale]
Sia $A\subset \RR^n$ aperto 
e sia $ f\colon A \to \RR^m$ una funzione 
tale che per ogni $ x\in A$ esistono 
tutte le derivate parziali $\partial f_k/\partial x_j$
(la funzione è \emph{derivabile}).
Supponiamo inoltre che le derivate parziali siano 
tutte continue.
Allora $ f$ è differenziabile in ogni punto. 
\end{theorem}
\begin{proof}
\emph{Passo 1:} facciamo il caso $n=2$, $m=1$.
Scriveremo $(x,y)$ al posto di $ x$ 
e $(x_0,y_0)$ al posto di $ x_0$.
Il vettore obliquo $(x,y)-(x_0,y_0)$ 
può essere decomposto nella somma dei 
due vettori paralleli agli assi coordinati
$(x,y)-(x_0,y)$ e $(x_0,y)-(x_0,y_0)$.
Alla funzione ristretta 
alle direzioni orizzontale e verticale possiamo
applicare il teorema di Lagrange per ottenere
\begin{align*}
  f(x,y) - f(x_0,y_0)
  = f(x,y) - f(x_0,y) + f(x_0,y) - f(x_0,y_0) \\
  = f_x(\tilde x,y)\cdot (x-x_0)
    + f_y(x_0,\tilde y)\cdot (y-y_0)
\end{align*}
con $\tilde x \in (x_0,x)$ e $\tilde y\in(y_0,y)$.
I valori di $\tilde x$ e $\tilde y$ dipendono 
da $x$ e da $y$, ma tendono rispettivamente 
a $x_0$ e $y_0$ quando $(x,y)\to (x_0,y_0)$.
Vogliamo dimostrare che $f$ è differenziabile in $(x_0,y_0)$.
Se lo è, necessariamente il suo differenziale
si scriverà nella forma $df[h,k] = h f_x + k f_y$.
La formula~\eqref{eq:def_differenziale}
diventa quindi per $(x,y)\to(x_0,y_0)$
\begin{align*}
  \lefteqn{f(x,y) - f(x_0,y_0) 
  - f_x(x_0,y_0)\cdot (x-x_0) - f_y(x_0,y_0)\cdot (y-y_0)}\\
  &= \Enclose{f_x(x_0,y_0) - f_x(\tilde x, y)}(x-x_0)
   + \Enclose{f_y(x_0,y_0) - f_y(x_0,\tilde y)}(y-y_0)\\
  & = o(1) \cdot (x-x_0) + o(1) \cdot (y-y_0)
\end{align*}
in quanto per l'ipotesi di continuità delle derivate parziali
le parentesi quadre tendono a zero quando $(x,y)\to (x_0,y_0)$.
Osservando infine che 
$\abs{x-x_0}, \abs{y-y_0}\le \sqrt{(x-x_0)^2 + (y-y_0)^2}$
si ottiene il risultato voluto.

\emph{Passo 2:} $n>2$, $m=1$.
Si procede esattamente come nel passo 1, salvo che per 
passare dal punto $ x_0$ al punto $ x$ bisogna 
seguire $n$ segmenti paralleli agli assi su ognuno 
dei quali si applica il teorema di Lagrange.

\emph{Conclusione.} Se $m>1$ basta applicare il risultato 
alle componenti $f_1,\dots, f_m$ di $ f$ e osservare 
che la formula \eqref{eq:def_differenziale} è equivalente 
alle $m$ formule:
\[
  f_k( x) = f_k( x_0) + (d f_k)_{ x_0}[ x- x_0] + o(\abs{ x- x_0}).
\]
\end{proof}

\begin{example}
Si consideri la funzione $f(x,y) = (x-1)(y+2)$.
Allora si ha
\[
f_x(x,y) = \frac{d}{dx} f(x,y) = y + 2, 
\quad f_y(x,y) = \frac{d}{dy} f(x,y) = x - 1.
\]
Visto che le derivate parziali sono continue la funzione 
risulta essere differenziabile e $\nabla f(x,y) = (y+2,x-1)$.
\end{example}

\begin{example}
Si consideri la funzione $f(x,y) = (x-1)(y+2)$ 
e la curva $u(t) = (t, t^2)$.
Si ha 
\[
 \frac{d}{dt}f(u(t)) 
  = \Enclose{(t-1)(t^2+2)}'
  = (t^3 - t^2 + 2t - 2)'
  = 3t^2 - 2t + 2.
\]
Tralasciando le dipendenze si scriverà più semplicemente
$f = (x-1)(y+2)$, $x=t$, $y=t^2$.
La formula per la derivata della funzione composta
diventa quindi
\begin{align*}
  \frac{df}{dt} 
  &= \frac{\partial f}{\partial x}\cdot \frac{dx}{dt}
  + \frac{\partial f}{\partial y}\cdot \frac{dy}{dt} \\
  &= (y+2) + (x-1) 2t = t^2 + 2 + (t - 1)2t \\
  &= 3t^2 - 2t + 2. 
\end{align*}

Visto che $x=t$ possiamo anche più direttamente scrivere $y=x^2$.
In tal caso si scriverà
\begin{align*}
  \frac{df}{dx} 
  &= \frac{\partial f}{\partial x}
  + \frac{\partial f}{\partial y}\cdot \frac{dy}{dx} \\
  &= (y+2) + (x-1) 2x = x^2 + 2 + (x - 1)2x \\
  &= 3x^2 - 2x + 2.
\end{align*}
e si capisce quindi la necessità di distinguere tra derivata parziale 
$\frac{\partial f}{\partial x}$ e derivata totale $\frac{df}{dx}$.
\end{example}

\section{derivate successive}

In generale possiamo definire il differenziale di una funzione 
$f\colon V\to W$ definita tra due spazi vettoriali normati $V$ e $W$.
Il differenziale risulta essere allora 
una funzione $df\colon V \to \L(V,W)$ che manda 
ogni punto $x\in V$ in cui la funzione è differenziabile
nell'applicazione lineare e continua $Df(x) \in \L(V,W)$.
Lo spazio $\L(V,W)$ risulta a sua volta essere uno spazio normato
(dalla norma operatoriale). 
E' quindi possibile iterativamente differenziare il differenziale.

Nel caso in cui $W = \RR$ si definisce $V ' = \L(V,\RR)$ (lo spazio duale) 
e dunque il differenziale $Df$ è a sua volta una funzione $Df\colon V \to V'$. 
Lo spazio duale può essere messo in corrispondenza con $V$ se 
rappresentiamo ogni $L \in V'$ mediante un vettore $\xi \in V$ 
tramite il prodotto scalare su $V$: $L[v] = \langle \xi, v\rangle$.
Il vettore di $V$ che rappresenta il differenziale è il gradiente $\nabla f$:
$Df{x}[v] = \langle \nabla f, v\rangle$.
Il differenziale del differenziale in ogni punto è una funzione lineare $D^2 f(x)\in \L(V,V')$
ovvero $D^2 f \colon V \to \L(V,V')$. 
Un elemento di $\L(V,V')$ rappresenta una forma bilineare 
perché ad ogni $v \in V$ viene associata una applicazione lineare 
che ad ogni $w\in V$ restituisce uno scalare:
\[
  D^2 f(x)[v,w] 
  = \langle D^2 f(x)[v], w\rangle
  = \frac{\partial^2 f}{\partial w \partial v}(x).
\]
Dunque $D^2f(x)$ si rappresenta mediante una matrice 
chiamata \emph{matrice hessiana} le cui componenti 
sono
\[
  \enclose{D^2 f(x)}_{jk}
  = D^2 f(x)[e_j,e_k] = \frac{\partial^2 f}{\partial x_k \partial x_j}
\]
ovvero
\[
  D^2 f = \begin{pmatrix}
    \frac{\partial^2 f}{(\partial x_1)^2}, \dots, \frac{\partial^2 f}{\partial x_n \dots \partial x_1} \\
    \vdots, \ddots, \vdots \\
    \frac{\partial^2 f}{\partial x_n \partial x_1}, \dots, \frac{\partial^2 f}{(\partial x_n)^2}
  \end{pmatrix}.
\]

Per le derivate $k$-esime il differenziale $D^k f$ rappresenta
una forma $k$-lineare $(v_1, \dots, v_k) \mapsto D^k f(x)[v_1, \dots, v_k]$
e si rappresenta mediante un tensore a $k$ indici.

\begin{definition}[classi di regolarità]
Sia $A\subset \RR^n$ un aperto. 
Definiamo $C^0(A,\RR^m)$ l'insieme delle funzioni continue
$ f\colon A \to \RR^m$. 
Per induzione si definisce la classe $C^k(A,\RR^m)$ 
delle funzioni che hanno tutte le prime $k$ derivate parziali 
continue:
\[
   C^k(A,\RR^m) = \ENCLOSE{ f\colon A\to\RR^m\colon
      f_{x_1}, \dots,  f_{x_n} \in C^{k-1}(A,\RR^m)
   }.
\]
Definiamo inoltre
\[
  C^\infty(A, \RR^m) = \ENCLOSE{ f\colon A\to \RR^m\colon 
  f\in C^k(A,\RR^m)\,\forall k\in \NN}.
\]

Si intende poi $C^k(A) = C^k(A,\RR)$, $C^\infty(A) = C^\infty(A,\RR)$.
\end{definition}

\begin{theorem}[Schwarz]
Sia $A\subset \RR^n$ un aperto 
e sia $f\colon A \subset \RR^n \to \RR$.
Fissati $k,j\in \ENCLOSE{1,\dots,n}$
supponiamo che in tutti i punti di $A$ esistano le derivate parziali 
$\frac{\partial f}{\partial x_k}$ 
e $\frac{\partial f}{\partial x_j}$
e che a loro volta esistano le derivate seconde miste 
$\frac{\partial^2 f}{\partial x_j \partial x_k}$
e $\frac{\partial^2 f}{\partial x_k \partial x_j}$.
Supponiamo infine che le derivate seconde miste siano continue 
in un punto $ x_0\in A$. Allora sono uguali:
\[
 \frac{\partial^2 f}{\partial x_j \partial x_k}( x_0)
 = \frac{\partial^2 f}{\partial x_k \partial x_j}( x_0).
\]
\end{theorem}
%
\begin{proof}
  \emph{Passo 1:} supponiamo $n=2$ e usiamo le variabili $x,y$ invece che 
  $x_1,x_2$. 
L'idea è che le derivate seconde miste approssimano l'incremento che si ottiene 
per andare dal punto $(x_0,y_0)$ al punto $(x,y)$ seguendo due diversi percorsi,
il primo incrementando prima la $x$ e poi la $y$ e il secondo 
incrementando prima la $y$ e poi la $x$.
Consideriamo quindi la seguente identità:
\[
\scriptstyle
\Enclose{f(x,y)-f(x_0,y_0)} - \Enclose{f(x,y_0) - f(x_0,y_0)}
= \Enclose{f(x,y) - f(x,y_0)} - \Enclose{f(x_0,y) - f(x_0, y_0)}
\]
e applichiamo il teorema di Lagrange 
a sinistra alla funzione $f(x,\cdot) - f(x_0,\cdot)$ nell'intervallo $[y_0,y]$ 
e a destra alla funzione $f(\cdot,y) - f(\cdot,y_0)$ nell'intervallo $[x_0,x]$.
Si ottiene quindi
\[
\Enclose{\frac{\partial f}{\partial y}(x,\tilde y)-\frac{\partial f}{\partial y}(x_0,\tilde y)}(y-y_0)
= \Enclose{\frac{\partial f}{\partial x}(\tilde x,y)-\frac{\partial f}{\partial x}(\tilde x,y_0)}(x-x_0)
\]
per qualche $\tilde x \in (x_0,x)$ e $\tilde y \in (y_0,y)$.
Applichiamo ora nuovamente il teorema di Lagrange a sinistra
sull'intervallo $[x_0,x]$ e a destra sull'intervallo $[y_0,y]$ 
per ottenere
\[
\frac{\partial^2 f}{\partial x\partial y}(\hat x,\tilde y)(x-x_0)(y-y_0)
= \frac{\partial^2 f}{\partial y\partial x}(\tilde x,\hat y)(y-y_0)(x-x_0)
\]
ovvero
\[
\frac{\partial^2 f}{\partial x\partial y}(\hat x,\tilde y)
= \frac{\partial^2 f}{\partial y\partial x}(\tilde x,\hat y)
\]
per qualche $\hat x \in (x_0,x)$ e $\hat y \in (y_0,y)$.

Per $(x,y)\to(x_0,y_0)$ si ha $\tilde x, \hat x \to x_0$ e $\tilde y, \hat y\to y_0$
e grazie all'ipotesi di continuità delle derivate seconde miste 
passando al limite si conclude:
\[
\frac{\partial^2 f}{\partial x\partial y}(x_0, y_0)
= \frac{\partial^2 f}{\partial y\partial x}(x_0, y_0).  
\]

\emph{Caso generale.}
Consideriamo la funzione di due variabili ottenuta da $f$ fissando 
tutte le coordinate diverse da $j$ e $k$ ai valori delle corrispondenti 
coordinate del punto $ x_0$. 
Si può quindi applicare il passo precedente.
\end{proof}

\begin{example}
Si consideri la funzione $f(x,y,z) = (x+y^2) \sin z$.
Si ha
\[
 f_x = \sin z, \qquad
 f_y = 2 y \sin z, \qquad
 f_z = (x+y^2)\cos z
\]
e
\begin{gather*}
  f_{xx} = 0, \qquad
  f_{xy} = 0, \qquad 
  f_{xz} = \cos z \\
  f_{yx} = 0, \qquad 
  f_{yy} = 2 \sin z, \qquad
  f_{yz} = 2y \cos z\\
  f_{zx} = \cos z, \qquad
  f_{zy} = 2y\cos z, \qquad
  f_{zz} = -(x+y^2)\sin z.
\end{gather*}
Dunque abbiamo verificato che le derivate 
seconde commutano: $f_{xy}=f_{yz}$, $f_{xz}=f_{zx}$,
$f_{yz}=f_{zy}$
e la matrice hessiana è simmetrica:
\[
 D^2 f(x,y) = \begin{pmatrix}
  0 & 0 & \cos z\\
  0 & 2\sin z & 2y \cos z \\
  \cos z & 2y \cos z & -(x+y^2)\sin z
 \end{pmatrix}.
\]
\end{example}

\begin{theorem}[Formula di Taylor]
Sia $A\subset \RR^n$ aperto e 
sia $f\colon A\subset \RR^n \to \RR$ una funzione di classe $C^2$.
Allora fissato $ x_0\in A$ si ha, per $ x \to  x_0$
\begin{equation}\label{eq:taylor_secondo}
  f( x) = f( x_0) + \langle \nabla f( x_0),  x -  x_0\rangle 
  + \frac 1 2 \langle D^2 f( x_0) ( x -  x_0),  x -  x_0\rangle
  + o(\abs{ x -  x_0}^2).
\end{equation}

Più in generale, se $f$ è di classe $C^m$ si avrà
\[
 f(x) = f(x_0) + \sum_{k=1}^m \frac 1 {k!} D^k f(x_0)[x- x_0, \dots, x-x_0]
  + o(\abs{ x -  x_0}^k).
\]
\end{theorem}
%
\begin{proof}
\emph{Caso $m=2$ (secondo ordine)}.
Fissato $ x$ (oltre che $ x_0$) consideriamo 
il vettore $ v =  x -  x_0$ e restringiamo la funzione 
alla retta passante per $ x_0$ con direzione $ v$.
Otteniamo una funzione 
\[
  g(t) = f( x_0 + t  v)
\]
le cui derivate sono:
\begin{align*}
  g'(t) &= \langle \nabla f( x_0 + t v),  v\rangle,\\
  g''(t) &= \langle D^2 f( x_0+t v)[v],  v \rangle.
\end{align*}
Usando la formula di Taylor con resto di Lagrange per le funzioni di una variabile vista nel 
corso di Analisi Matematica Uno, si ottiene
\[
  g(t) = g(0) + g'(0) t + \frac 1 2 g''(\bar t( x)) t^2
\]
per un qualche $\bar t( x)\in (0,t)$. 
Abbiamo quindi 
\[
  f( x) = f( x_0) + \langle \nabla f(x_0),  x -  x_0\rangle 
  + \frac 1 2 \langle D^2 f( x_0 + \bar t( x)(x -  x_0))[x- x_0],  x- x_0\rangle.
\]
Visto che per $ x \to  x_0$ si ha $\bar t( x)\to 0$,
per la continuità delle derivate seconde
si ha 
\[
  D^2 f( x_0 + \bar t( x)( x -  x_0))
  \to D^2 f( x_0)
\]
e dunque segue la tesi.

\emph{Caso generale.}
La dimostrazione è analoga, con una discreta complicazione delle notazioni.
\end{proof}

\section{ottimizzazione}

Data una funzione $f \colon A\subset \RR^n \to \RR$ ci poniamo 
il problema di determinare eventuali punti di massimo e minimo di $f$ 
su $A$.

Se $x_0\in A$ diremo che $x_0\in A$ è un minimo assoluto (o globale) 
di $f$ su $A$ se $f(x_0) \le f(x)$ per ogni $x\in A$. 
Diremo che è un massimo assoluto se $f(x_0)\ge f(x)$ per ogni $x\in A$.
Diremo che $x_0$ è un punto di massimo (o minimo) relativo (o locale)
se è un massimo (o minimo) assoluto per $f$ in un intorno di $x_0$.

\begin{theorem}[condizioni al primo ordine]
Sia $f\colon A\subset \RR^n \to \RR$ e sia $x_0$ un punto 
di massimo o minimo relativo di $f$.
Se $x_0$ è un punto interno di $A$ e se $f$ è derivabile in $x_0$ 
allora $\nabla f(x_0)=0$.
Un punto $x_0$ in cui si ha $\nabla f(x_0) = 0$ si chiama 
\emph{punto critico}.
\end{theorem}
\begin{proof}
Se $x_0$ è minimo (o massimo) relativo per $f$
allora la funzione $g_k(t) = f(x_0+te_k)$ ha minimo (o massimo)
relativo nel punto $t=0$. 
Tale funzione è definita in un intorno di $t=0$ dunque 
per i criteri già visti nel corso di Analisi Matematica Uno
si deve avere $g_k'(0)=0$. 
Ma $g'_k(0) = \frac{\partial f}{\partial x_k}(x_0)$
Questo è vero per ogni $k=1,\dots,n$, si ottiene quindi $\nabla f(x_0)=0$.
\end{proof}

